\section{Background}
\label{sec:bg-full}

\subsection{Branching in Agentic and Reasoning Workloads}
\label{sec:bg-search}

Chain-of-thought produces one trace~\cite{cot}; Tree-of-Thoughts, self-consistency, and Monte Carlo tree search produce many and then select~\cite{tot,selfconsistency,mctsllm}.
In the Game-of-24 setting of Tree-of-Thoughts, a proposal prompt emits candidate equations, a value prompt labels each candidate, and breadth-first search retains a beam~\cite{tot}.
Agent harnesses follow the same pattern: a coding or tool-using agent proposes several continuations of one context, executes or observes one, and discards the rest~\cite{react,sweagent,openhands}.
On GSM8K and MATH the accuracy gain of such width is well documented~\cite{gsm8k,math,math500}, and so is its token cost~\cite{equivpruner,stoploss}.

\subsection{Decode-Time Path Pruning}
\label{sec:bg-prune}

Early-stopping self-consistency draws a window of samples and stops once the window agrees~\cite{esc,selfconsistency}.
Speculative Rejection ranks partial generations at a reward horizon and halts the lower quantile~\cite{specrej}.
DPTS generates a mini-step on every active node and early-stops low-confidence children~\cite{dpts}.
STOP organizes these signals by source and learnability and trains an internal token that predicts futility from a short prefix~\cite{stoploss}; EquivPruner merges semantically equivalent actions~\cite{equivpruner}.
Each signal is a function of generated text.
On our GSM8K mix at $k{=}16$, ESC, Speculative Rejection, and DPTS therefore prefill the same $5008$ tokens and differ only in decode volume ($131072$, $98304$, and $78336$ tokens).

\subsection{Prefix Caching and Disaggregated Transfer}
\label{sec:bg-pd}

APC and RadixAttention index full blocks by chained keys, $h_j=\mathrm{hash}(h_{j-1},\text{block}_j,\text{extra})$~\cite{vllm,sglang}.
A request reuses a block only if the block and its entire parent chain are resident.
Whether $k$ siblings issued together share the trunk depends on the cache regime of \S\ref{sec:regimes}: a resident trunk (R3) is shared by all siblings; a cold trunk is computed once if siblings are serialized behind it (R2), and between one and $k$ times if they are admitted concurrently, depending on how the scheduler treats blocks that are still being computed (R1).
Prefix caching thus addresses trunk reuse, but it neither rejects a residual nor decides which pages cross a connector.

PD disaggregation places prefill and decode on separate instances and transfers KV between them~\cite{distserve,splitwise}.
Every admitted request is prefilled on $\Pin$, and the pages that $\Din$ lacks are transferred; a residual that the harness later discards has already consumed both.
Program- and workflow-level schedulers raise the scheduling unit above the request~\cite{autellix,thunderagent,continuum,mori}, and KV retention removes trunk recomputation across tool gaps; the remaining per-step cost is the residual suffixes of siblings that the harness can enumerate before the observation returns.

\subsection{Prefill, Decode, and Idle Time}
Prefill processes the prompt in parallel and is compute-bound; decode emits one token per step and is memory-bandwidth-bound~\cite{orca,distserve,splitwise}.
PagedAttention stores KV in non-contiguous pages~\cite{vllm}, and chunked prefill co-schedules a prompt chunk with a decode batch~\cite{sarathi}.
For agents, the prompt of step $t{+}1$ extends that of step $t$ after a tool interval of variable length~\cite{continuum,mori}, during which known suffixes can be prefilled (Appendix~\ref{sec:schedule}).

\subsection{Where Existing Mechanisms Act}
\label{sec:traditional-app}

\Cref{fig:trad} places each mechanism on the stage timeline of one fan-out: the cost incurred before the earliest point at which a non-winner can be discarded.

\begin{figure*}[t]
\centering
\resizebox{\textwidth}{!}{% Where each method pays for one fan-out: first expand / rejection / connector.
% Visual replacement for the deleted main-paper Table 1.
\begin{tikzpicture}[
  x=1cm, y=1cm,
  font=\sffamily,
  >=Latex,
  line cap=round, line join=round,
  note/.style={font=\tiny\sffamily, align=center},
  chip/.style={draw, rounded corners=1pt, inner sep=1.4pt, font=\tiny\sffamily,
              align=center, line width=0.4pt},
  rowlab/.style={font=\tiny\sffamily\bfseries, anchor=east, align=right},
  stage/.style={font=\tiny\sffamily, text=fsgray},
]
\path (0,0) rectangle (16.45,7.35);

% title strip
\node[font=\scriptsize\bfseries\sffamily, anchor=west] at (0.06,7.15)
  {Where the fan-out is paid};
\node[note, text=fsgray, anchor=east] at (16.38,7.15)
  {one fan-out of $k$ residuals on trunk $L$};

% stage axis
\draw[fsgray!40, line width=0.5pt] (3.20,6.55) -- (15.90,6.55);
\foreach \x/\lab in {
  4.10/clone $L$,
  6.40/prefill $\rho$,
  8.70/\texttt{insert},
  11.00/generate,
  13.50/reject} {
  \fill[fsgray!55] (\x,6.55) circle (0.07);
  \node[stage] at (\x,6.85) {\lab};
}
\node[note, text=fsgray, anchor=west] at (0.10,6.55) {stage $\rightarrow$};

% legend
\node[chip, draw=fsmiss!80!black, fill=fsmiss] at (4.55,6.15) {paid};
\node[chip, draw=fsorange!70, fill=fsspec] at (6.35,6.15) {reject here};
\node[chip, draw=fsgreen!55!black, fill=fscommit] at (8.55,6.15) {connector};
\node[chip, draw=fsblue!50, fill=fslight] at (11.05,6.15) {\sys};

% helper: draw a row
% args via foreach below

% --- Recompute ---
\node[rowlab, text=fsgray] at (3.05,5.45) {Recompute};
\draw[fsgray!25, rounded corners=1.5pt, fill=black!2] (3.20,5.10) rectangle (15.90,5.80);
\fill[fsmiss] (3.40,5.28) rectangle (5.00,5.62);
\node[font=\tiny\sffamily, text=fsred] at (4.20,5.45) {$k{\times}L$};
\fill[fsmiss] (5.20,5.28) rectangle (7.60,5.62);
\node[font=\tiny\sffamily, text=fsred] at (6.40,5.45) {$k$ residuals};
\node[chip, draw=fsorange!70, fill=fsspec] at (13.50,5.45) {end of request};
\node[note, text=fsgray, anchor=west] at (14.55,5.45) {no connector};

% --- APC / radix ---
\node[rowlab, text=fsblue] at (3.05,4.45) {APC / radix};
\draw[fsblue!15, rounded corners=1.5pt, fill=fslight] (3.20,4.10) rectangle (15.90,4.80);
\fill[fsmiss] (3.40,4.28) rectangle (5.00,4.62);
\node[font=\tiny\sffamily, text=fsred] at (4.20,4.45) {$k{\times}L$};
\fill[fsmiss] (5.20,4.28) rectangle (7.60,4.62);
\node[font=\tiny\sffamily, text=fsred] at (6.40,4.45) {first expand};
\node[chip, draw=fsorange!70, fill=fsspec] at (11.00,4.45) {LRU after publish};
\node[note, text=fsblue, anchor=west] at (12.55,4.45) {replay of $\star$ hits};

% --- P/D disagg ---
\node[rowlab, text=fspurple] at (3.05,3.45) {P/D disagg.};
\draw[fspurple!12, rounded corners=1.5pt, fill=white] (3.20,3.10) rectangle (15.90,3.80);
\fill[fsmiss] (3.40,3.28) rectangle (5.00,3.62);
\node[font=\tiny\sffamily, text=fsred] at (4.20,3.45) {$k{\times}L$};
\fill[fsmiss] (5.20,3.28) rectangle (7.60,3.62);
\node[font=\tiny\sffamily, text=fsred] at (6.40,3.45) {prefill all};
\fill[fsgreen!45] (7.80,3.28) rectangle (9.60,3.62);
\node[font=\tiny\sffamily, text=fsgreen!25!black] at (8.70,3.45) {$kL{+}\sum\ell_i$};
\node[chip, draw=fsorange!70, fill=fsspec] at (11.00,3.45) {after insert};
\node[note, text=fspurple, anchor=west] at (12.55,3.45) {aborts still paid};

% --- ESC / SR / DPTS ---
\node[rowlab, text=fsorange!85!black] at (3.05,2.45) {ESC / SR / DPTS};
\draw[fsorange!12, rounded corners=1.5pt, fill=fsspec!40] (3.20,2.10) rectangle (15.90,2.80);
\fill[fsmiss] (5.20,2.28) rectangle (7.60,2.62);
\node[font=\tiny\sffamily, text=fsred] at (6.40,2.45) {every $\rho_i$};
\fill[fsgreen!45] (7.80,2.28) rectangle (9.60,2.62);
\node[font=\tiny\sffamily, text=fsgreen!25!black] at (8.70,2.45) {KV of all};
\fill[fsorange!35] (9.80,2.28) rectangle (12.20,2.62);
\node[font=\tiny\sffamily, text=fsorange!85!black] at (11.00,2.45) {prefix};
\node[chip, draw=fsorange!70, fill=fsspec] at (13.50,2.45) {after generate};

% --- PD-Prune ---
\node[rowlab, text=fsgreen!35!black] at (3.05,1.30) {\sys};
\draw[fsgreen!25, rounded corners=1.5pt, fill=fscommit, line width=0.55pt] (3.20,0.70) rectangle (15.90,1.90);
\fill[fsblue!55] (3.40,1.30) rectangle (5.00,1.70);
\node[font=\tiny\sffamily, text=white] at (4.20,1.50) {alias $L$};
\fill[fsgreen!55] (5.20,1.30) rectangle (7.20,1.70);
\node[font=\tiny\sffamily, text=white] at (6.20,1.50) {miss of $\Keep$};
\fill[fsgreen!40] (7.40,1.30) rectangle (9.40,1.70);
\node[font=\tiny\sffamily, text=fsgreen!20!black] at (8.40,1.50) {$\Xfer$ of $\Keep$};
\node[chip, draw=fsgray, fill=black!4, text=fsgray] at (6.40,0.95) {$\Lambda$ before kernel};
\node[chip, draw=fsorange!70, fill=fsspec] at (11.80,1.50) {$t_0$ probe};
\node[note, text=fsgreen!30!black, anchor=west] at (13.00,1.50) {radix + decode prune};

% footer
\node[note, text=fsgray, anchor=west] at (0.10,0.30)
  {Red = first-expand cost.\quad
   Orange chip = earliest rejection.\quad
   Green bar = connector payload.\quad
   \sys row: \S\ref{sec:app}.};
\end{tikzpicture}
}
\caption{\textbf{Where each mechanism acts on one fan-out of $k$ residuals over a trunk of length $L$.}
Red bars denote prefill issued before any sibling is discarded, orange chips the earliest rejection point, and green bars connector volume.
Trunk volumes for the prefix-cache rows assume the cold concurrent regime R1; in regimes R2--R4 the trunk is computed at most once (\S\ref{sec:regimes}).
The \sys{} row follows \S\ref{sec:app}.}
\label{fig:trad}
\end{figure*}

\subsection{Notation}
\label{sec:model}

\begin{table}[t]
\centering
\caption{Notation.}
\label{tab:notation}
\scriptsize
\setlength{\tabcolsep}{3.2pt}
\begin{tabular}{@{}llp{0.52\columnwidth}@{}}
\toprule
Symbol & Domain & Meaning \\
\midrule
$L$, $\ell_i$, $\ell_\star$ & $\mathbb{N}$ & trunk; residual $i$; committed residual \\
$x_u$, $\rho_i$, $x_\star$ & tokens & trunk; residual; spine \\
$\Pin$, $\Din$ & --- & prefill and decode instances \\
$c_u^{\mathcal{X}}$, $m_i^{\mathcal{X}}$ & tokens & resident prefix of trunk and residual on $\mathcal{X}$ \\
$s_i$, $s'_i$ & $[0,1]$ & residual score; probe score \\
$\tau_{\mathrm{pre}}$, $\tau_{\mathrm{dec}}$ & $[0,1]$ & rejection and extension thresholds \\
$a_i$ & $\Act$ & \Hit, \Admit, \Probe, \Reject (\Cref{def:action}) \\
$e_i$ & $\{0,1\}$ & probe of $i$ is extended \\
$t_0$, $D$, $d_i$ & tokens & probe length; budget; decode tokens of $i$ \\
$W_i$, $X_i$ & tokens & computed prefill; connector tokens \\
$W_i^{\mathrm{ext}}$ & tokens & recomputation on probe extension \\
$\Rej$, $\mathcal{H}$, $\mathcal{E}$ & $\subseteq[k]$ & rejected, hit, and extended residuals \\
$b$, $P$ & bytes, tokens & KV bytes per token; page size \\
$M_{\mathrm{pool}}$ & bytes & KV pool capacity \\
$B_t=B^c_t+B^s_t$ & tokens & committed and speculative chunk budget \\
\bottomrule
\end{tabular}
\end{table}
